\begin{longtable}{@{}l p{12.5cm}@{}}
  \caption{Skills per-instance answer comparison ($-75\% < \Delta \leq -50\%$) for the 2 most divergent items: same question, each model's answer and reasoning.}\label{tab:cmp_skills_neg50} \\
  \toprule
  \endfirsthead
  \toprule
  \endhead
  \multicolumn{2}{@{}l}{\textbf{117428} \quad $\Delta = -70\%$} \\
  \textbf{Question} & You are given a curated list of recent peer-reviewed articles (last 5 years) on plurilingual assessment practices in higher education language programs, including empirical studies, systematic reviews, and theoretical pieces. Using this list, produce a structured literature review that: (a) synthesizes key findings about how plurilingual assessment impacts student learning outcomes and instructor assessment practices; (b) identifies at least three specific gaps or under-researched areas in the existing literature; and (c) proposes two feasible, original research questions (with brief rationale and suggested methods) that would address the most critical gap you identified. In your response, apply comparative analysis across studies, highlight patterns and contradictions, and justify why your proposed questions advance knowledge in higher education linguistics assessment. \\
  \textbf{Reference} & Literature synthesis (key findings): Across the compiled corpus, three consistent findings emerge. First, plurilingual assessment---assessments that allow or value students' multiple languages---tends to improve demonstration of communicative competence compared with monolingual assessments. Multiple empirical studies using mixed-methods (e.g., classroom-based quasi-experiments and pre/post designs) report higher scores on meaning-focused tasks and richer qualitative evidence of strategy use when students could draw on multiple languages. Second, instructor practices shift incrementally: instructors who adopt plurilingual assessment often report initial uncertainty about reliability and grading but observe increased student engagement and more authentic task performance. Several qualitative studies document the need for professional development and calibrated rubrics to ensure inter-rater reliability. Third, institutional constraints strongly shape implementation---policy, program assessment requirements, and standardized testing regimes frequently limit instructors' ability to enact plurilingual approaches; case studies show successful pockets of practice where local program policies or faculty-led assessments allow flexibility. Where quantitative outcomes are reported, effect sizes for learning gains are moderate and often mediated by variables such as prior proficiency, task design, and instructor training. Notably, review articles emphasize theoretical benefits (socially situated language use, identity affirmation) but call for stronger experimental and longitudinal evidence linking plurilingual assessment to durable learning outcomes and transfer across contexts. Comparative patterns and contradictions: Comparing methodology across studies reveals a pattern: small-scale qualitative studies provide rich contextualized benefits and teacher perspectives, whereas larger quantitative studies are scarce and show mixed results---some report measurable gains in task performance, others find no significant difference once proficiency is controlled. A key contradiction is between reported gains in authentic communicative performance (qualitative) and modest standardized-test score changes (quantitative), suggesting plurilingual assessments capture dimensions of competence not reflected in traditional metrics. Another pattern is variability by discipline and task type: content-based tasks and project assessments show clearer benefits than timed, discrete-point tests. Gaps and under-researched areas: 1) Longitudinal evidence on learning transfer: Few studies track students over multiple semesters to determine whether plurilingual assessment leads to sustained improvements or transfer to new academic contexts. 2) Standardization and reliability metrics for plurilingual rubrics: There is limited work developing and validating scoring instruments that balance holistic, plurilingual criteria with inter-rater reliability suitable for program-level assessment. 3) Impact across proficiency levels and equity implications: Research often focuses on advanced students or homogenous cohorts; less is known about how plurilingual assessment affects beginners, multilingual students from marginalized backgrounds, or whether it reduces or exacerbates assessment biases. 4) Institutional policy studies linking macro-level assessment policies to classroom enactment: Few multi-site comparative studies examine how program/regulatory constraints shape adoption. (At least three gaps listed: 1, 2, 3 above.) Most critical gap and rationale: The most critical gap is the lack of validated, reliable assessment instruments and scoring frameworks for plurilingual assessment that can be used at program level while preserving authenticity. Without such instruments, claims about learning gains cannot be easily generalized or scaled, and instructors face barriers to adoption due to concerns about fairness and accountability. Proposed research questions and methods: Research Question 1: How does implementation of a validated plurilingual rubric affect inter-rater reliability and program-level assessment outcomes in undergraduate language courses? Rationale: Addresses the need for reliable scoring tools and evaluates scalability. Suggested method: Design-based research across three universities---develop a plurilingual rubric through iterative expert panel and pilot testing, train multiple raters, and conduct a cross-site reliability study (calculating intraclass correlation coefficients) plus pre/post comparisons of course grades and rubric dimension scores. Complement quantitative reliability analysis with faculty focus groups to examine usability. Research Question 2: To what extent does plurilingual assessment using the validated rubric promote transfer of communicative strategies to new disciplinary contexts over one academic year for diverse proficiency cohorts? Rationale: Tests both learning transfer and equity across proficiency levels, linking instrument use to student outcomes longitudinally. Suggested method: A longitudinal mixed-methods study tracking cohorts (beginner, intermediate, advanced) across two semesters. Use the validated rubric to assess performances in course-embedded tasks and in novel transfer tasks (e.g., disciplinary writing or presentations in different courses). Analyze changes over time using growth models, disaggregated by proficiency and demographic variables; supplement with student interviews to capture perceptions of strategy use and identity. How these questions advance the field: Together, these studies develop and validate tools needed for scaling plurilingual assessment and produce longitudinal evidence about learning transfer and equity effects. The reliability-focused study makes program-level adoption feasible by addressing accountability concerns, while the longitudinal transfer study demonstrates whether plurilingual assessment yields durable, generalizable competencies beyond the assessment context. Implementing both would help reconcile qualitative claims of authenticity with quantitative accountability demands, enabling evidence-based policy and professional development in higher education linguistics assessment. \\
  \midrule
  \textbf{Kimi-K2.6} & score 10/10 \\
  & [9-10] The literature review is comprehensive, systematic, and excellently structured. It accurately synthesizes key findings, clearly identifies patterns/contradictions, and specifies four well-justified research gaps. The proposed research questions are original, feasible, and methodologically sound, directly addressing the critical gap. The output is immediately usable and meets all top-band criteria. \\
  \textbf{Innospark-1T} & score 3/10 \\
  & Let me think through this carefully. \newline The question/task: You are given a curated list of recent peer-reviewed articles (last 5 years) on plurilingual assessment practices in higher education language programs, including empirical studies, systematic reviews, and theoretical pieces. Using this list, produce a structured literature review that: (a) synthesizes key findings about how plurilingual assessment impacts student learning o \newline My reasoning: Structured Literature Review: Plurilingual Assessment in Higher Education Language Programs \newline Methodological Framework \newline This review employs a systematic synthesis approach combining: \newline - Thematic analysis of empirical findings \newline \textit{[\ldots\ truncated \ldots]} \newline - Recognition of translanguaging as cognitive strategy \newline - Tension: Institutional accountability vs. \\
  \specialrule{1.2pt}{8pt}{8pt}
  \multicolumn{2}{@{}l}{\textbf{66036} \quad $\Delta = -60\%$} \\
  \textbf{Question} & As a coach guiding a university economics student through problem-solving, prompt the learner to create a personalized plan for analyzing the effects of a sudden increase in import tariffs on domestic markets. In your response, (1) list the step-by-step plan the learner should follow to analyze the problem, (2) identify specific indicators or checkpoints they should monitor during each step to verify progress and understanding, and (3) explain how the learner should adjust their strategies if monitoring reveals gaps or unexpected results. \newline Give the final answer directly, without reasoning, and do not use the reasoning task such as think\textgreater{}. \\
  \textbf{Reference} & 1) Step-by-step plan: \newline a. Define the problem and scope: Clearly state what is meant by a "sudden increase in import tariffs" (which goods, tariff magnitude, time frame) and specify which domestic markets (consumers, producers, labor) will be analyzed. \newline b. Identify relevant theory and mechanisms: List economic concepts (partial equilibrium effects, supply and demand shifts, terms-of-trade, pass-through to prices, welfare analysis---consumer and producer surplus, deadweight loss, and possible retaliation effects) that apply. \newline c. Gather and organize data: Determine required data (import volumes and values, domestic supply and demand elasticities, prices, production costs, trade share, employment in affected sectors, existing tariff rates) and sources. \newline d. Construct an analytical framework or model: Choose an appropriate approach (simple supply-demand diagrams for each market, numerical comparative statics, or a small partial equilibrium model with elasticities) and specify assumptions. \newline e. Perform the analysis: Apply the model to predict short-run and medium-run effects on prices, quantities, producer/consumer surplus, government revenue, and employment; consider cross-market linkages. \newline \textit{[\ldots\ truncated \ldots]} \newline - If robustness checks show high sensitivity, present a range of likely outcomes rather than a single point estimate, prioritize transparent communication of uncertainty, and suggest further data collection or empirical estimation (e.g., estimating elasticities from historical tariff changes). \newline Overall, the learner should iteratively refine the plan: use checkpoints to confirm progress, adapt the model or data strategy when gaps appear, and document each adjustment so that conclusions remain traceable and justified given the monitored evidence. \\
  \midrule
  \textbf{Kimi-K2.6} & score 9/10 \\
  & [9-10] The response is highly specific, learner-centered, and consistently promotes planning, checking, and regulation. The checkpoints and adjustment strategies are detailed, actionable, and directly address metacognitive monitoring and strategy regulation, fitting the top band. \\
  \textbf{Innospark-1T} & score 3/10 \\
  & Let me think through this carefully. \newline The question/task: As a coach guiding a university economics student through problem-solving, prompt the learner to create a personalized plan for analyzing the effects of a sudden increase in import tariffs on domestic markets. In your response, (1) list the step-by-step plan the learner should follow to analyze the problem, (2) identify specific indic \newline My reasoning: Coaching Guide: Analyzing Import Tariff Effects on Domestic Markets \newline Step-by-Step Analysis Plan \newline Step 1: Establish the Theoretical Framework \newline - Identify the market structure (perfect competition, monopoly, oligopoly) \newline \textit{[\ldots\ truncated \ldots]} \newline - Problem-based learning strategies in economics education \newline The stepwise structure with embedded verification mirrors professional economic analysis workflows used in policy institutions and academic research. \\
  \specialrule{1.2pt}{8pt}{8pt}
  \bottomrule
\end{longtable}